\documentclass[11pt,a4paper]{article}
\usepackage[T1]{fontenc}
\usepackage{lmodern}
\usepackage[margin=25mm,headheight=14pt]{geometry}
\usepackage{amsmath,amssymb,amsthm,mathtools,booktabs,array,longtable}
\usepackage{microtype,xcolor,enumitem,fancyhdr,xurl}
\usepackage[colorlinks=true,linkcolor=blue!45!black,urlcolor=blue!55!black,citecolor=blue!45!black]{hyperref}
\hypersetup{pdftitle={A 4 by 4 Obstruction to Polynomial Ordered Binary Matrix Removal},pdfauthor={Research investigation, 9 October 2026},pdfsubject={Complete mathematical proof, prior-work audit, and reproducible exact certificates}}
\setlength{\parskip}{3pt}
\setlength{\parindent}{1.2em}
\setlist{itemsep=3pt,topsep=5pt}
\pagestyle{fancy}\fancyhf{}
\fancyhead[L]{\small Ordered binary matrix removal}
\fancyhead[R]{\small Research record, 9 October 2026}
\fancyfoot[C]{\thepage}
\newtheorem{theorem}{Theorem}[section]
\newtheorem{lemma}[theorem]{Lemma}
\newtheorem{proposition}[theorem]{Proposition}
\newtheorem{corollary}[theorem]{Corollary}
\theoremstyle{remark}
\newtheorem{remark}[theorem]{Remark}
\newcommand{\one}{\mathbf{1}}
\newcommand{\eps}{\varepsilon}
\newcommand{\dist}{\operatorname{dist}}
\newcommand{\Rrole}{\mathsf R}
\newcommand{\Crole}{\mathsf C}
\newcommand{\ancR}{\widehat R}
\newcommand{\ancC}{\widehat C}
\newcommand{\crown}{\mathcal C}
\newcommand{\Hfour}{H_4}
\title{\textbf{A \(4\times4\) Obstruction to Polynomial\\Ordered Binary Matrix Removal}\\[6pt]
\large Sparse zero localization despite unpinned anchors}
\author{Research record prepared in this investigation}
\date{9 October 2026}
\begin{document}
\maketitle
\begin{abstract}
The strongest prior result located in this investigation is an October 7
ProveIt research manuscript giving a \(5\times5\) obstruction, improving
OpenAI's \(66\times66\) construction. That manuscript also displays a
\(4\times4\) candidate and an unpinned copy that survives its proposed
repair. We prove that this previously displayed candidate does fail
polynomial ordered binary matrix removal in the original two-group host.
The new structural fact is that every copy, including every unpinned
copy, maps its prescribed zero at coordinate \((3,4)\) to a leaf row and
its parent-minus column block. The order of the two minus-parity modes
excludes the only possible nonterminal exception.

For \(h\ge1\), put \(m=2^h\), \(d_h=14h+2\), and \(n_h=d_hm\).
The explicit hosts have normalized edit distance at least \(d_h^{-2}\)
from \(\Hfour\)-freeness under arbitrary binary changes. Every copy
meets a fixed \(2m\)-cell set at a specified pattern coordinate, so
\(N_{\Hfour}(A_h)\le2mn_h^6\). A complete three-family classification
further gives \(N_{\Hfour}(A_h)=(41m^6-74m^5)/4\) for \(h\ge2\),
and \(N_{\Hfour}(A_1)=128\).

All proofs are included. Two separately implemented exhaustive programs
agree on the complete weighted copy sets through depth five. Their
finite checks support the implementations; they do not replace the
all-depth proof. The new theorem has not been kernel-checked in Lean.
No minimality or exhaustive publication-priority claim is made.
\end{abstract}
\newpage
{\small\tableofcontents}
\newpage

\section{Strongest baseline, contribution, and precise statement}

\subsection{Dated prior-work finding}
As checked on 9 October 2026, current \texttt{openai/math} main is
\path{fd4aeeb2ee4fc729c18d98444fed42fd0529eeeb}, the snapshot supplied
for this investigation. Its family-190 manuscript and Lean catalogue
still specify a \(66\times66\) counterexample~\cite{oai,oaiLean}.
The public main history has three commits and one merged pull request.
The family-190 paths have not changed since the initial public commit.

A stronger external result was found through GitHub code search:
\emph{A \(5\times5\) Obstruction to Polynomial Ordered Matrix Removal},
a ProveIt research draft dated 7 October 2026~\cite{proveit}.
The first returned path-history commit is
\texttt{bfbfe915186310fdf745a7a089e7dd88eeb773b1},
dated 7 October at 21:24:26 UTC. It supplies ordinary proofs for
\[
H_5=
\begin{pmatrix}
0&1&1&1&0\\
1&0&1&0&1\\
1&1&0&0&0\\
1&0&0&1&0\\
0&1&0&1&1
\end{pmatrix},
\qquad n=(20h+2)2^h,
\]
with exactly \(2(2^h)^8\) copies and minimum repair size \(4^h\).
Its main structural tool is full-pattern pinning with the crown
\(J_3-I_3\). We reconstructed its principal pinning, counting, and
arbitrary-edit proofs and found no gap in those statements. Its code was
also rerun successfully; that rerun is distinguished from our independent
implementations. The prior result is not Lean-formalized.

The same prior manuscript already displays the matrix \(\Hfour\) below,
and gives an unpinned copy surviving a dummy-row-by-leaf-column repair.
It leaves the possibility of a different counting argument open.
\textbf{The candidate matrix itself is prior work. The contribution here
is proving that this candidate works, by classifying and controlling its
unpinned copies.}

The first concrete dense-anchor reduction considered was a six-position
dense prefix with 32 binary-word rows, giving a \(38\times38\) anchor.
That reduction is valid, as detailed in Appendix~\ref{app:38}, but is
superseded by the located \(5\times5\) theorem. The precise stronger size
target therefore became \(4\times4\).

\subsection{Definitions and every quantifier}
Write \([n]=\{1,\ldots,n\}\). For a binary \(k\times k\) matrix \(H\)
and a binary \(n\times n\) matrix \(A\), an ordered copy is a pair of
strictly increasing maps \(r,c:[k]\to[n]\) satisfying
\[
A(r(i),c(j))=H(i,j)\quad\text{for every }i,j\in[k].
\]
The maps are chosen independently. Every zero and every one must match.
Rows and columns need not be consecutive. There is no quotient by
permutations or other symmetries. Let \(N_H(A)\) count these pairs.
The matrix \(A\) is \(H\)-free exactly when \(N_H(A)=0\).

For \(n\ge1\), define
\[
\dist_H(A)=\frac1{n^2}
 \min_{\substack{B\in\{0,1\}^{[n]\times[n]}\\N_H(B)=0}}
 \bigl|\{(x,y)\in[n]^2:A(x,y)\ne B(x,y)\}\bigr|.
\]
For our patterns the minimum is nonempty because the all-zero matrix is
free. Edits can be anywhere and can change bits in either direction.
The normalization of copy counts is \(n^{2k}\); using
\(\binom nk^2\) instead only changes fixed-pattern constants in the
relevant range.

A fixed \(H\) has polynomial ordered removal if
\[
\begin{split}
\exists c_H,C_H>0\ \forall n\in\mathbb N_{\ge1}\
\forall\eps\in(0,1)\
\forall A\in\{0,1\}^{[n]\times[n]}:\qquad\\
\dist_H(A)\ge\eps\ \Longrightarrow\
N_H(A)\ge c_H\eps^{C_H}n^{2k}.
\end{split}
\]
Failure for the explicit \(\Hfour\) means
\[
\begin{split}
\forall c,C\in\mathbb R_{>0}\ \exists n\in\mathbb N_{\ge1}\
\exists\eps\in(0,1)\ \exists A\in\{0,1\}^{[n]\times[n]}:\qquad\\
\dist_{\Hfour}(A)\ge\eps
\quad\text{and}\quad N_{\Hfour}(A)<c\eps^C n^8.
\end{split}
\]
The constructed sequence in fact supplies such witnesses with arbitrarily
large \(n\) and arbitrarily small positive \(\eps\).

\begin{theorem}[Four-by-four obstruction]\label{thm:main}
Let
\[
\boxed{\Hfour=
\begin{pmatrix}
0&1&1&0\\
1&0&0&1\\
1&0&1&0\\
0&1&1&1
\end{pmatrix}.}
\]
For every integer \(h\ge1\), set
\[
m=2^h,\qquad d_h=14h+2,\qquad n_h=d_hm,\qquad \eps_h=d_h^{-2}.
\]
The deterministic host \(A_h\) in Section~\ref{sec:host} satisfies
\[
\dist_{\Hfour}(A_h)\ge\eps_h,\qquad
N_{\Hfour}(A_h)\le 2m n_h^6.
\]
More precisely,
\[
N_{\Hfour}(A_h)=
\begin{cases}
128,&h=1,\\[2pt]
\dfrac{41m^6-74m^5}{4},&h\ge2.
\end{cases}
\]
Consequently
\[
\frac{N_{\Hfour}(A_h)}{n_h^8}
\le \frac{41}{4}\eps_h^4\,4^{-h},
\]
and polynomial ordered binary matrix removal fails for \(\Hfour\).
\end{theorem}

\section{Reconstruction of the published baseline}
\label{sec:baseline}

\subsection{The 64-by-64 anchor and 2-by-2 body}
For zero-based row and column indices \(u,v\in\{0,\ldots,63\}\), the
published anchor is \(J_{64}-I_{64}\), except that rows \(32,\ldots,63\)
and columns \(59,\ldots,63\) are replaced by all five-bit words:
\[
S_{64}(u,59+j)=
 \left\lfloor\frac{u-32}{2^{4-j}}\right\rfloor\bmod2
\quad(32\le u<64,\ 0\le j<5).
\]
With \(e_1,e_2\) the first two coordinate vectors, the pattern is
\[
H_{66}=
\begin{pmatrix}
S_{64}&e_1&e_2\\
e_1^{\mathsf T}&1&0\\
e_2^{\mathsf T}&1&1
\end{pmatrix},\qquad
P=\begin{pmatrix}1&0\\1&1\end{pmatrix}.
\]
Rows and columns of \(S_{64}\) are distinct. Each row has at least
58 ones, each column at least 48. The first 32 positions on either
axis give a dense core: any row or column has at most one zero there.
The 32 word rows shatter the final five columns.

\subsection{Host, distance, and count are separate ingredients}
The host uses the signed binary-tree positions and six modes per level
given explicitly below. For a general anchor order \(s\), there are
\(6hs\) anchor groups, \(2h+1\) variable classes, and one dummy class
per axis, each of size \(m=2^h\). Thus
\[
n=(6hs+2h+2)m.
\]
For \(s=64\), this is \((386h+2)m\).

The published distance proof selects a common random path and independent
anchor, dummy, and within-block representatives. Every protected changed
cell is encountered with probability \(m^{-2}\). Fewer than \(m^2\)
edits therefore leave one protected selection intact. Avoiding its four
types of \(P\)-body would propagate opposite root values to the same
leaf cell, a contradiction. This proof applies to the edited matrix
itself, including copies that edits create.

The published counting proof is different. A nonanchor variable position
has at most four incidences with selected distinct anchor groups.
Nonanchor rows do not shatter five nonanchor columns. The dense core and
word block therefore pin every leading \(S_{64}\)-copy to one complete
mode anchor. Its body then has the designated roles. Three mode types
are impossible; the fourth forces a common diagonal leaf cell.
There are \(m\) such cells, and hence
\[
N_{H_{66}}(A_h)\le mn^{130},\qquad
\frac{N_{H_{66}}(A_h)}{n^{132}}
\le (386h+2)^{-2}2^{-h}.
\]
Changing the original leaf-diagonal ones to zeros hits all original
copies but creates plus-mode copies. A transversal is not a repair.

This distinction is consistent with Alon--Ben-Eliezer's Theorem 1.2,
which concerns entry-disjoint existing copies, whereas their Problem 1.4
asks about genuine edit distance~\cite{ab}. Qualitative ordered removal
is established in~\cite{abf}. The singleton polynomial question is
Conjecture 4.5 in the first version of~\cite{gs}. The ordered-graph
classification in~\cite{gt} concerns one common vertex order, a different
model from independent row and column selections.

\subsection{What the actual Lean implementation proves}
The solution module
\path{OAI.Combinatorics.MatrixRemoval.Main} proves
\path{OAI.Problem348.counterexample_sequence} and
\path{OAI.Problem348.no_polynomial_removal_bound}.
The underlying definitions use two strictly increasing maps and Boolean
equality at every cell. Hamming distance counts all differing cells.
The minimum ranges over all binary matrices, assigning nonfree matrices
the sentinel \(n^2+1\); the all-zero free matrix makes this equivalent
to the stated minimum when \(n\ge1\).

The sequence theorem is unconditional for the fixed \(66\times66\)
pattern. Its concrete host, all-copy bound, and arbitrary-repair bound
are discharged by the imported implementation. The comparator file's
\texttt{sorry} is a challenge placeholder, not the solution proof.
Static inspection of 61 MatrixRemoval implementation files, totaling
7,609 lines, found no \texttt{sorry}, \texttt{admit}, \texttt{axiom}, or
\texttt{native\_decide}. The comparator permits
\texttt{propext}, \texttt{Quot.sound}, and \texttt{Classical.choice};
it has \texttt{enable\_nanoda=false}. The toolchain is Lean 4.34.1.
\textbf{No Lean build or axiom-printer invocation was executed here.}
These are source-level findings, not a newly rerun kernel certificate.

\section{The explicit two-group host}\label{sec:host}

\subsection{Positions, multiplicities, and order}
Fix \(h\ge1\) and \(m=2^h\). At depth \(i\), nodes have indices
\(0\le p<2^i\); write \(\pi(p)=\lfloor p/2\rfloor\).
For \(0\le i<h\), each node has disjoint row blocks
\(R_i^-(p),R_i^+(p)\) and column blocks \(C_i^-(p),C_i^+(p)\),
all of size \(m/2^i\). At depth \(h\), the signs are identified:
\[
R_h^-(p)=R_h^+(p)=R_h(p),\qquad
C_h^-(p)=C_h^+(p)=C_h(p),
\]
and each leaf block has size one. These are the only identifications.

A variable class \(R_i^\sigma\) or \(C_i^\sigma\) is the union of
all blocks at a fixed depth and sign. The shared leaf class occurs once.
Every class has size \(m\). For \(i\ge1\), a subscript \(d\in\{0,1\}\)
restricts to nodes \(p\equiv d\pmod2\). Add dummy row and column blocks
\(R_*,C_*\), each of size \(m\).

The recursive row order at a nonleaf node is its minus block, the
zero-child subtree, the one-child subtree, then its plus block.
The column order is its plus block, the same two subtrees, then its
minus block. Place a leaf block once. Put \(R_*\) after all variable
rows and \(C_*\) before all variable columns. Within each block use
indices \(0,\ldots,\lvert\text{block}\rvert-1\) in increasing order.

For \(1\le i\le h\), comparisons with a parent-level block give
\begin{equation}\label{eq:orders}
\begin{aligned}
R_i^+(p)<R_{i-1}^+(a)&\iff\pi(p)\le a,&
R_{i-1}^-(a)<R_i^-(p)&\iff a\le\pi(p),\\
C_{i-1}^+(b)<C_i^+(q)&\iff b\le\pi(q),&
C_i^-(q)<C_{i-1}^-(b)&\iff\pi(q)\le b.
\end{aligned}
\end{equation}
Here \(X<Y\) means every position of the first block precedes every
position of the second. The formulas hold for shared leaves too.
Indeed, depth-\((i-1)\) subtrees occur in increasing node order, and
each child lies between its parent's opening and closing blocks.
At a fixed nonleaf depth \(j\),
\begin{equation}\label{eq:same-level}
C_j^-(a)<C_j^+(b)\iff a<b,\qquad
C_j^+(a)<C_j^-(b)\iff a\le b.
\end{equation}
Complements of the comparisons in \eqref{eq:orders} give the strict
reverse inequalities used below; the compared blocks are disjoint.

\subsection{Modes and anchor groups}
At every depth \(i=1,\ldots,h\), introduce the six modes, in this order:
\[
V_i^+,\ V_i^-,\ W_{i,0}^+,\ W_{i,1}^+,\ W_{i,0}^-,\ W_{i,1}^-.
\]
Modes at smaller depths come first. Use this same mode order on both
axes. Their four role sets are
\begin{equation}\label{eq:roles}
\begin{array}{c|cc|cc}
t&\Rrole_1(t)&\Rrole_2(t)&\Crole_1(t)&\Crole_2(t)\\ \hline
V_i^+&R_i^+&R_{i-1}^+&C_*&C_{i-1}^+\\
V_i^-&R_{i-1}^-&R_i^-&C_*&C_{i-1}^-\\
W_{i,d}^+&R_{i,d}^+&R_*&C_{i-1}^+&C_{i,d}^+\\
W_{i,d}^-&R_{i,d}^-&R_*&C_{i,d}^-&C_{i-1}^-
\end{array}
\end{equation}
For each mode, the two row roles are disjoint, as are the two column
roles. Correct ordering of representatives is a separate requirement;
whole role sets need not precede one another.

Before every variable and dummy position, place two anchor groups per
mode on each axis, denoted \(\ancR_t(1),\ancR_t(2)\) and
\(\ancC_t(1),\ancC_t(2)\). Every group has size \(m\); group 1 precedes
group 2. Thus there are \(12h\) anchor groups and \(2h+2\)
nonanchor classes per axis:
\[
n=(12h+2h+2)m=(14h+2)m.
\]
Identify the ordered axes with \([n]\).

\subsection{Every entry}
Between anchors set
\[
A_h(\ancR_t(u),\ancC_{t'}(v))
=\one[t=t']\one[u\ne v].
\]
A nonanchor row in \(\Rrole_j(t)\) has signature \(e_j^{\mathsf T}\)
against the two anchor columns of \(t\); outside both roles it has
signature zero. A nonanchor column in \(\Crole_j(t)\) has signature
\(e_j\) against its anchor rows; otherwise its signature is zero.

Between nonanchors put one whenever at least one position is dummy.
All remaining nonzero entries are specified in this table:
\begin{equation}\label{eq:entries}
\begin{array}{ccc|c}
\text{row block}&\text{column block}&\text{range}&A_h\\ \hline
R_i^+(p)&C_i^+(q)&0\le i\le h&\one[p\le q]\\
R_i^-(p)&C_i^-(q)&0\le i<h&\one[p<q]\\
R_i^+(p)&C_{i-1}^+(q)&1\le i\le h&\one[\pi(p)\le q]\\
R_i^-(p)&C_{i-1}^-(q)&1\le i\le h&\one[\pi(p)<q]
\end{array}
\end{equation}
Cases absent from the table are zero. Entries are constant on a pair
of blocks. The same-level minus line excludes the leaves, so the
shared leaf-by-leaf threshold is weak. In the last two lines the
column depth is below \(h\), and its two signs are distinct. Thus the
assignments never conflict.

\section{Structural facts valid at every depth}\label{sec:structural}

\begin{lemma}[Twins, nesting, and signatures]\label{lem:structure}
The following hold for every \(h\ge1\).
\begin{enumerate}
\item Positions within any one anchor group are full row or column
twins. Positions in a dummy group are full twins too.
\item An anchor position's positive nonanchor neighborhood is contained
in one variable class, a parity subclass, or the dummy class.
\item Within a variable row class, traces on opposite nonanchors are
nested by inclusion as the node index increases. The same statement
holds for variable columns. Within any variable class there are at
most two full traces on opposite anchors.
\item For a fixed \emph{nonanchor} row, its entries along a variable
column class are nondecreasing in node order. A nonleaf variable row
can be nonzero only in its own signed column class and its signed
parent class. A leaf row can meet only the leaf class and the two
signed parent classes.
\item A row in \(\Rrole_2(t)\) and a column in \(\Crole_1(t)\) always
have entry one.
\end{enumerate}
\end{lemma}
\begin{proof}
The first two statements follow from group-constant entries and the role
table. The variable thresholds in \eqref{eq:entries} are nonincreasing
in the row node and nondecreasing in the column node. Their zero
patterns depend only on the class; dummy incidences are constant.
This proves nesting and the nonanchor-row monotonicity assertion.
Anchor rows themselves need not have monotone variable traces, because
their incidences can select parity. The stated restriction is essential.

For completeness, the variable mode incidences, with undefined levels
omitted, are as follows; \((t,j)\) means role \(j\):
\[
\begin{array}{c|l}
R_i^+(p),\,i<h&(V_i^+,1),(V_{i+1}^+,2),(W_{i,d}^+,1)\\
R_i^-(p),\,i<h&(V_i^-,2),(V_{i+1}^-,1),(W_{i,d}^-,1)\\
R_h(p)&(V_h^+,1),(V_h^-,2),(W_{h,d}^+,1),(W_{h,d}^-,1)
\end{array}
\]
and
\[
\begin{array}{c|l}
C_i^+(p),\,i<h&(V_{i+1}^+,2),(W_{i+1,0}^+,1),
 (W_{i+1,1}^+,1),(W_{i,d}^+,2)\\
C_i^-(p),\,i<h&(V_{i+1}^-,2),(W_{i+1,0}^-,2),
 (W_{i+1,1}^-,2),(W_{i,d}^-,1)\\
C_h(p)&(W_{h,d}^+,2),(W_{h,d}^-,1).
\end{array}
\]
Here \(d=p\bmod2\). Only parity varies within a fixed class, proving
the two-signature bound, including at the shared leaves. These lists
also verify the four-mode incidence bound used in the original dense
anchor argument. Finally, in a \(V\)-mode column role 1 is dummy, and
in a \(W\)-mode row role 2 is dummy, proving the last assertion.
\end{proof}

\begin{lemma}[Nonanchor crown exclusion]\label{lem:crown}
The nonanchor submatrix contains no row-and-column permutation of
\(\crown_3=J_3-I_3\).
\end{lemma}
\begin{proof}
Such a crown would have three distinct two-element row traces on three
columns. A dummy column is impossible because it is one on every
nonanchor row. Two columns from the same variable class are impossible
because their traces are nested, whereas a crown has incomparable
column traces. Thus the columns occupy three distinct variable classes.
A dummy row has trace \(111\), so cannot supply a crown row.

Join two column classes whenever the support of some variable row can
meet both. This graph is \emph{contained in} the graph formed by two
signed paths
\[
C_0^\sigma-C_1^\sigma-\cdots-C_{h-1}^\sigma
\quad(\sigma=+,-)
\]
and the terminal triangle
\[
T=\{C_{h-1}^+,C_{h-1}^-,C_h\}.
\]
The containing graph has only this triangle. Some displayed edges
need not be realized at small depths; equality of graphs is not used.
The three required two-element supports would force the selected
classes to be \(T\). Every variable row able to meet two classes of
\(T\) is a leaf row. Its trace is a permutation of
\[
\bigl(\one[\pi(p)\le q_+],\one[\pi(p)<q_-],\one[p\le q_0]\bigr),
\]
which is nested as \(p\) increases. It cannot supply three incomparable
two-element traces. This proves the lemma, including \(h=1\).
\end{proof}

\section{Every ordered copy localizes to a sparse zero set}
\label{sec:localization}

\begin{theorem}[All-depth sparse zero localization]\label{thm:localization}
In every ordered \(\Hfour\)-copy in \(A_h\), its third row is a leaf
row \(R_h(p)\), and its fourth column belongs to
\(C_{h-1}^-(\pi(p))\). Thus the prescribed zero at pattern coordinate
\((3,4)\) lies in
\[
Z_h=\bigcup_{p=0}^{m-1}R_h(p)\times C_{h-1}^-(\pi(p)),
\qquad |Z_h|=2m.
\]
\end{theorem}

\subsection{Exhaust all possible anchor counts}
\begin{lemma}\label{lem:shape}
Every copy has exactly two anchor rows and either one or two anchor
columns. With two anchor columns, both axes use the two complete groups
of one mode and the other positions have that mode's designated roles.
\end{lemma}
\begin{proof}
Let \(a,b\) be the total numbers of anchors among its four rows and four
columns. They are prefix counts because all anchors come first.
All four rows of \(\Hfour\) are distinct, as are all four columns.

Suppose \(a\le1\). If \(b=0\), rows \(2,3,4\) and columns \(1,3,4\)
are nonanchors and form
\[
\begin{pmatrix}1&0&1\\1&1&0\\0&1&1\end{pmatrix},
\]
a permuted crown, contradicting Lemma~\ref{lem:crown}. If
\(1\le b\le2\), the first anchor column has ones in nonanchor rows
2 and 3, putting them in one nonanchor class. On nonanchor columns
3 and 4 their traces are \(01\) and \(10\), impossible for nested
variable traces or dummy twins. If \(b\ge3\), anchor column 1 connects
rows 2 and 3, and anchor column 3 connects rows 3 and 4. The three
nonanchor rows consequently lie in one class. Their traces on the
first three anchor columns are \(100,101,011\), three distinct
signatures, again impossible. Hence \(a\ge2\).

When \(a,b\ge2\), let \(t_1\le t_2\) be the first two row modes and
\(u_1\le u_2\) the first two column modes. The ones at \((1,2)\)
and \((2,1)\) give \(t_1=u_2\) and \(t_2=u_1\). Common mode
order forces all four to agree. Distinct full traces prohibit reuse
of an anchor group; order therefore gives groups 1 and 2 on both
axes. An extra anchor body row with a unit signature would repeat
one of those groups and thus a full row trace. An anchor in another
mode has zero signature. Neither can be a body row. The identical
argument applies to columns, so \(a=b=2\), and signatures identify
the four body roles. This step uses full twins, not just the leading
two-by-two anchor.

It remains to consider \(b\le1\). If \(b=0,a=2\), anchor row 2 puts
nonanchor columns 1 and 4 in one class. Their traces on nonanchor
rows 3 and 4 are \(10,01\), impossible. If \(b=0,a\ge3\), the
positive supports of the first three anchor rows are
\(\{2,3\},\{1,4\},\{1,3\}\). Their overlaps put all four columns
in one class, but their traces on these three anchor rows are
\(011,100,101,010\), four distinct signatures. This is impossible
for either a variable or dummy class.

Finally, if \(b=1,a\ge3\), the ones at \((2,1)\) and \((3,1)\)
put anchor rows 2 and 3 in the same complementary group of the first
anchor column's mode. They would be full twins, contrary to the
pattern. The only noncanonical possibility is \((a,b)=(2,1)\).
\end{proof}

\subsection{Canonical copies: two anchor columns}
For a canonical copy, the usual four-mode calculation applies with
no lower bound on the anchor size. In \(V_i^+\), let \(p\) be the
child row node, \(a\) the parent row node, and \(b\) the parent column
node. Row order and the bottom-right one give
\(\pi(p)\le a\le b\), so the top-right entry is one, impossible.
In \(V_i^-\), row order and the top-right zero give
\(b\le a\le\pi(p)\); the bottom-right strict threshold is zero,
also impossible.

In \(W_{i,d}^+\), child row and column nodes \(p,q\) have the same
parity. The top-left one and column order give
\(\pi(p)\le b\le\pi(q)\), so \(p\le q\), contradicting the
top-right zero. In \(W_{i,d}^-\), the inequalities instead give
\(\pi(q)\le b\le\pi(p)\), hence \(p\ge q\). At \(i<h\) the
top-left entry \(\one[p<q]\) is zero. At \(i=h\) the weak leaf
threshold forces \(p=q\) and \(b=\pi(p)\). This establishes
Theorem~\ref{thm:localization} for canonical copies.

\subsection{Unpinned copies: one anchor column}
Write \(x,y\) for selected rows 3 and 4 and \(c_2,c_3,c_4\) for
the nonanchor columns. The one at \((2,1)\) puts anchor row 2 and
anchor column 1 in complementary groups of some mode \(t\).
If column 1 were group 2, then \(x\) would be in row role 2 and
\(c_4\) in column role 1 of \(t\). Their entry is always one by
Lemma~\ref{lem:structure}, contrary to \(H_4(3,4)=0\).
Consequently column 1 is group 1, row 2 is group 2, and
\begin{equation}\label{eq:unpinned}
x\in\Rrole_1(t),\quad c_4\in\Crole_2(t),\quad A_h(x,c_4)=0.
\end{equation}

The first anchor row has ones at \(c_2,c_3\), putting them in one
nonanchor class \(K\). Their entries against \(x\) are \(0,1\).
Thus \(K\) is variable, and \(x\) is not dummy. Also \(c_4\)
is not dummy since its entry against \(x\) is zero. Class \(K\)
differs from the class of \(c_4\): otherwise \(c_3<c_4\) and
nonanchor-row monotonicity would contradict their values \(1,0\).
Let \(t_0\) be the first anchor row's mode. Its order gives
\(t_0\le t\).

\paragraph{Case \(t=V_i^+\).}
Row \(x\) is child-plus node \(p\), possibly a leaf, and \(c_4\)
is parent-plus node \(b\). Its zero gives \(\pi(p)>b\).
If \(K\) is the child class, write \(q\) for \(c_3\)'s node.
Its one gives \(p\le q\), and \(c_3<c_4\) gives
\(\pi(q)<b\). Thus \(\pi(p)\le\pi(q)<b\), a contradiction.
Only at a leaf can \(K\) alternatively be the negative parent
class. Then its one gives \(\pi(p)<q\), while the same-level
minus-before-plus order gives \(q<b\). This also contradicts
\(\pi(p)>b\).

\paragraph{Case \(t=V_i^-\).}
Row \(x\) is a nonleaf minus node \(p\) at depth \(j=i-1\).
Column \(c_4\) is in its own minus class at node \(b\); its zero
gives \(p\ge b\). The only other possible support class is the
previous minus class, which does not exist when \(j=0\).
Otherwise let its node be \(q\). Its one and its order before
\(c_4\) give
\[
\pi(p)<q<\pi(b)\le\pi(p),
\]
a contradiction.

\paragraph{Case \(t=W_{i,d}^+\).}
Row \(x\) and column \(c_4\) are child-plus nodes \(p,b\) of the
same parity, possibly shared leaves. Their zero gives \(p>b\).
If \(K\) is the positive parent class, its node \(q\) satisfies
\[
\pi(p)\le q\le\pi(b).
\]
Equal parity would imply \(p\le b\), impossible. At a leaf the
negative parent class is also possible in principle; it instead
gives \(\pi(p)<q<\pi(b)\), again impossible.

\paragraph{Case \(t=W_{i,d}^-\), nonterminal \(i<h\).}
Row \(x\) is minus node \(p\) of parity \(d\), and \(c_4\) is
parent-minus node \(b\). Its zero gives \(b\le\pi(p)\).
The remaining class \(K\) must be the own minus class. If its
node at \(c_3\) is \(q\), the one and column order give
\[
p<q,\qquad \pi(q)\le b\le\pi(p).
\]
Parent monotonicity forces \(\pi(p)=\pi(q)=b\). Binary siblings
with \(p<q\) must be
\begin{equation}\label{eq:siblings}
p=2b,\qquad q=2b+1.
\end{equation}
Thus \(d=0\).

Now use the anchor mode order. A positive anchor incidence with
the nonterminal column class \(C_i^-\) is possible only in group 1
of \(W_{i,e}^-\), with columns restricted to parity \(e\), or in
group 2 of \(V_{i+1}^-\) or \(W_{i+1,e}^-\). The next-level
modes come after \(t=W_{i,0}^-\), so cannot supply the first anchor
row. Mode \(W_{i,1}^-\) also comes after \(t\). The first row
must therefore use \(e=0\), forcing \(q\) even, contrary to
\eqref{eq:siblings}. This excludes every nonterminal exception.

\paragraph{Case \(t=W_{h,d}^-\), terminal.}
Now \(x\) is leaf \(p\), and \(c_4\) is negative parent node \(b\)
with \(b\le\pi(p)\). If \(K\) is the leaf class, its node \(q\)
at \(c_3\) satisfies
\[
p\le q,\qquad \pi(q)\le b\le\pi(p),
\]
so \(b=\pi(p)=\pi(q)\). It is unnecessary, and generally false,
to force \(p=q\). If \(K\) is the positive parent class, its node
\(q\) gives
\[
\pi(p)\le q\le b\le\pi(p),
\]
so \(b=q=\pi(p)\). The negative parent class was already excluded
by monotonicity. These are every possible support class.

Both terminal possibilities put the fixed zero \((3,4)\) into
\(Z_h\), completing the proof of Theorem~\ref{thm:localization}.
Each leaf row has size one and its parent-minus column block has
size two, so \(|Z_h|=2m\).

\begin{corollary}[A sufficient copy bound]\label{cor:sparsecount}
For every \(h\ge1\), \(N_{\Hfour}(A_h)\le2mn^6\).
\end{corollary}
\begin{proof}
Fix the images of pattern row 3 and column 4 to be one of the
\(2m\) cells in \(Z_h\). The other three row images and three
column images have at most \(n^6\) choices. Ordering and matching
can only reduce this count.
\end{proof}

\section{Distance against arbitrary edits}\label{sec:distance}
\begin{theorem}\label{thm:distance}
For every \(h\ge1\), every \(\Hfour\)-free binary \(n\times n\)
matrix differs from \(A_h\) in at least \(m^2\) entries.
\end{theorem}
\begin{proof}
Let \(B\) be any binary matrix whose changed-cell set \(E\) has
\(|E|<m^2\). We prove \(B\) contains \(\Hfour\).

Partition each axis into its anchor groups, its dummy class, its
signed variable classes below the leaves, and its shared leaf class
counted once. Every class has size \(m\).
Choose a uniform leaf \(z\in\{0,\ldots,m-1\}\) and let
\[
a_i=\left\lfloor z/2^{h-i}\right\rfloor\quad(0\le i\le h).
\]
Conditionally on \(z\), choose uniform positions
\[
x_i^\sigma\in R_i^\sigma(a_i),\qquad
y_i^\sigma\in C_i^\sigma(a_i)
\]
independently within distinct blocks, using just one choice on each
axis at the shared leaf. Independently choose a uniform position
from every anchor group and both dummy blocks; call the latter
\(x_*,y_*\).

Each variable representative is uniform on its size-\(m\) class:
its node has probability \(2^{-i}\), and each within-block choice
has conditional probability \(2^i/m\). An anchor or dummy
representative is independent of every opposite-axis representative.
The root representatives are independent uniform choices in fixed
size-\(m\) blocks. Hence, for each fixed cell with at least one
anchor or dummy coordinate, or with both coordinates at root level,
\[
\mathbb P(\text{both coordinates are selected})=m^{-2}.
\]
Call these cells protected. Nonroot variable pairs may be correlated;
no \(m^{-2}\) claim is made for them.
There is one representative per class, including just one leaf
representative per axis. Summing over the set of changed protected
cells, each counted once, yields
\[
\mathbb P(\text{some selected protected cell belongs to }E)
\le |E|/m^2<1.
\]
Fix a selection avoiding every such change. Its anchor entries,
anchor-to-body signatures, dummy incidences, and root entries agree
with the original host. All remaining entries of \(B\) are unrestricted.

In each mode, the selected two anchor rows and columns form the
required \(\crown_2\), preceding the selected body positions.
Any correctly ordered selected body equal to
\(P=\left(\begin{smallmatrix}1&0\\1&1\end{smallmatrix}\right)\)
therefore extends to an actual \(\Hfour\)-copy in \(B\).
Suppose, toward a contradiction, that \(B\) is free.
Its unchanged roots give
\[
B(x_0^+,y_0^+)=1,\qquad B(x_0^-,y_0^-)=0.
\]
Assume these values have reached level \(i-1\).
The selected \(V_i^+\) and \(V_i^-\) bodies are
\[
\begin{pmatrix}1&B(x_i^+,y_{i-1}^+)\\1&1\end{pmatrix},
\qquad
\begin{pmatrix}1&0\\1&B(x_i^-,y_{i-1}^-)\end{pmatrix}.
\]
Their orders follow from the common parent-child path and
\eqref{eq:orders}; the left columns are intact dummy incidences.
Avoiding \(P\) forces the two unknown entries to be one and zero,
respectively. Put \(d=a_i\bmod2\). The selected \(W_{i,d}^+\)
and \(W_{i,d}^-\) bodies then are
\[
\begin{pmatrix}1&B(x_i^+,y_i^+)\\1&1\end{pmatrix},
\qquad
\begin{pmatrix}B(x_i^-,y_i^-)&0\\1&1\end{pmatrix}.
\]
Their bottom rows are intact dummy incidences, and their column
orders again follow from the common path. Avoiding \(P\) forces
\[
B(x_i^+,y_i^+)=1,\qquad B(x_i^-,y_i^-)=0.
\]
Induction reaches \(i=h\), where the two asserted values refer
to the same shared leaf cell. This is impossible. Thus \(B\)
contains a copy, proving the lower bound.
\end{proof}

This proof is independent of the all-copy classification in the
\emph{unedited} host. It allows arbitrary zero-to-one and one-to-zero
changes in every region. It never substitutes hitting original copies
for obtaining a free edited matrix. In particular, no claim is made
that flipping \(Z_h\), or the prior dummy-by-leaf rectangle, repairs
the two-group host. Nor do we claim the minimum repair size equals
\(m^2\); the lower bound is what the theorem needs.

\section{Complete copy classification and exact counts}
\label{sec:exact}

\begin{lemma}\label{lem:dummy}
The fourth selected row of every copy is dummy.
\end{lemma}
\begin{proof}
Canonical copies already have terminal minus-horizontal roles, so
the fourth row is dummy. In an unpinned terminal copy, \(x\) is
leaf row 3, and class \(K\) of columns 2 and 3 is either the leaf
class or the positive parent class. Row 4 has ones on \(K\) and
the negative parent column \(c_4\).
If \(K\) is the leaf class, any such nondummy row is a leaf, since
only leaf variable rows meet leaf columns. If \(K\) is the positive
parent class, any variable row meeting both positive and negative
parent classes must again be a leaf. A leaf row after \(x\) has
a larger node index; its entry at \(c_2\) is no greater than
\(x\)'s zero, contradicting the required one. Thus row 4 is dummy.
\end{proof}

Write \(M=m/2\). The following three families are exhaustive.
In each description choose arbitrary actual representatives of the
displayed blocks. The multiplicity of one block description is the
product of all eight block sizes. The full pattern has distinct
rows and columns, so no copy can select a twin block twice.

\subsection{Canonical family}
For every leaf \(p=2b+d\), use both anchor groups on both axes of
\(W_{h,d}^-\), then leaf row \(p\), a dummy row, leaf column \(p\),
and a column of \(C_{h-1}^-(b)\). Conversely these always form
a copy. There are \(m\) block descriptions, each of multiplicity
\(2m^5\), for
\begin{equation}\label{eq:canonicalcount}
2m^6
\end{equation}
expanded copies.

\subsection{Unpinned family with two positive-parent columns}
Choose
\[
0\le a<b<M,\qquad d\in\{0,1\},\qquad p=2b+d.
\]
The columns, in order, are group 1 of \(W_{h,d}^-\),
\(C_{h-1}^+(a)\), \(C_{h-1}^+(b)\), \(C_{h-1}^-(b)\).
The last three rows are group 2 of that minus mode, leaf \(p\),
and the dummy. The first row is precisely one of
\[
\ancR_{V_h^+}(2),\quad
\ancR_{W_{h,0}^+}(1),\quad
\ancR_{W_{h,1}^+}(1),
\]
or, when \(h\ge2\) and \(a\equiv b\equiv e\pmod2\),
\(\ancR_{W_{h-1,e}^+}(2)\).

To verify necessity, localization gives the parent index of \(c_4\)
as \(b=\pi(p)\). The zero and one on \(c_2,c_3\), together with
their order before \(c_4\), force their indices to be \(a<b\) and
\(b\), respectively. The displayed first rows are exactly the
anchor incidences with both positive-parent columns, as seen from
the complete incidence list. They precede the second row. The role
table checks every entry, proving sufficiency too.

Thus the number of block descriptions is
\[
6\binom M2+
2\left\{\binom{\lceil M/2\rceil}{2}
       +\binom{\lfloor M/2\rfloor}{2}\right\}.
\]
Each has multiplicity \(8m^4\). For \(h\ge2\), \(M\) is even,
and the expanded contribution is
\begin{equation}\label{eq:pluscount}
7m^6-16m^5.
\end{equation}
For \(h=1\), no pair \(a<b\) exists, so this family is empty.

\subsection{Unpinned family with two leaf columns}
Choose \(p=2b+d\), two leaf nodes \(q_0<p\le q_1\) with
\(\pi(q_1)=b\), and common column parity
\(q_0\equiv q_1\equiv e\pmod2\).
The columns are group 1 of \(W_{h,d}^-\), leaves \(q_0,q_1\),
and \(C_{h-1}^-(b)\). The last three rows again are group 2 of
that minus mode, leaf \(p\), and dummy. The first row is
\(\ancR_{W_{h,e}^+}(2)\), or, if \(e\le d\),
\(\ancR_{W_{h,e}^-}(1)\).

Necessity follows from the leaf threshold, localization, and the fact
that these are exactly the anchor incidences with leaf columns.
The inequality \(e\le d\) is precisely the first-anchor-row order
condition for a minus mode; all plus modes precede the minus modes.
The role and threshold tables directly verify sufficiency.
For each parent \(b\), the possibilities are
\[
\begin{array}{c|c|c}
p&q_1&\text{number of descriptions}\\ \hline
2b+1&2b+1&2b\\
2b&2b&2b\\
2b&2b+1&b
\end{array}
\]
because there are \(b\) earlier leaves of the required parity,
with two possible first rows in the first two cases and one in
the last. Summing gives \(5\binom M2\) descriptions, each of
multiplicity \(2m^4\), for
\begin{equation}\label{eq:leafcount}
\frac54m^6-\frac52m^5.
\end{equation}
This formula is zero at \(m=2\), as required.

\subsection{Conclusion and comparison of bounds}
The shape lemma, all-depth mode analysis, and dummy-row lemma prove
that the three descriptions are exhaustive. They are disjoint by
anchor count and the class of their two middle columns.
Adding \eqref{eq:canonicalcount}--\eqref{eq:leafcount} gives
\[
N_{\Hfour}(A_h)=\frac{41}{4}m^6-\frac{37}{2}m^5
\quad(h\ge2).
\]
At \(h=1\) only the canonical family remains, giving \(128\).
In particular \(N_{\Hfour}(A_h)\le(41/4)m^6\) for every \(h\ge1\).
Dividing by \(n^8=d_h^8m^8\) proves
\[
\frac{N_{\Hfour}(A_h)}{n^8}
\le\frac{41}{4}d_h^{-8}4^{-h}
=\frac{41}{4}\eps_h^4\,4^{-h}.
\]

The exact count is not necessary for the main size improvement.
Corollary~\ref{cor:sparsecount} already gives
\[
\frac{N_{\Hfour}(A_h)}{n^8}\le2\eps_h2^{-h}.
\]
For arbitrary \(c,C>0\), either bound divided by \(\eps_h^C\)
tends to zero: an exponential in \(h\) dominates every fixed
real power of \(14h+2\). Taking \(h\) sufficiently large proves
the fully quantified failure statement in Section 1, and completes
Theorem~\ref{thm:main}.

\begin{corollary}[Arbitrarily large orders at each fixed depth]
For every \(h\ge1\) and every integer \(t\ge1\), replace each entry
of \(A_h\) by a constant \(t\times t\) block, with consecutive
clones on each axis. The resulting matrix has order \(tn_h\),
normalized distance at least \(\eps_h\), and normalized copy
count exactly \(N_{\Hfour}(A_h)/n_h^8\).
\end{corollary}
\begin{proof}
All partition classes now have size \(tm\). In the protected-path
proof, choose the common tree leaf uniformly from the original
\(m\) leaves, then independent uniform representatives within
the enlarged blocks, including the size-\(t\) leaf blocks.
Every protected-cell selection probability is \((tm)^{-2}\).
The same propagation contradiction proves a lower bound
\((tm)^2\) on edits, and hence the same normalized distance.
Two clones of one original row or column cannot occur in an
\(\Hfour\)-copy, since its row and column traces are distinct.
Every copy therefore projects to a unique original copy, and
each original copy has exactly \(t^8\) lifts. Division by
\((tn_h)^8\) proves the count assertion.
\end{proof}
Thus the example also defeats formulations that allow an arbitrary
order threshold depending on \(\eps\): first choose \(h\) to make
the normalized copy density small, then choose \(t\) to exceed that
threshold. No interchange of \(\eps\)- and \(n\)-quantifiers is needed.

\section{Adversarial order control and what has not been proved}
\label{sec:audit}

\subsection{A negative control for the parity order}
The critical parity-order step is not cosmetic. Reverse
\(W_{i,0}^-\) and \(W_{i,1}^-\) at each level on both axes, leaving
every entry rule unchanged. For every \(h\ge3\), the following
fixed-depth positions form an ordered copy:
\[
\begin{array}{c|l}
\text{rows}&
\ancR_{W_{2,1}^-}(1),\
\ancR_{W_{2,0}^-}(2),\
R_2^-(2),\ R_*\\
\text{columns}&
\ancC_{W_{2,0}^-}(1),\
C_2^-(1),\ C_2^-(3),\ C_1^-(1).
\end{array}
\]
Their traces are \(0110,1001,1010,0111\).
All selected variable positions are nonleaves. The multiplicity is
\[
\left(m\cdot m\cdot\frac m4\cdot m\right)
\left(m\cdot\frac m4\cdot\frac m4\cdot\frac m2\right)
=\frac{m^8}{128}.
\]
Hence the altered construction has copy density at least
\(1/(128d_h^8)\), and cannot satisfy our superpolynomial copy
estimate. This proves failure of the localization mechanism under
that alteration; it does not prove polynomial removal for the
pattern in all other hosts.

\subsection{Weak steps examined}
\begin{enumerate}
\item \textbf{Unintended copies.} The anchor itself is not rigid.
Lemma~\ref{lem:shape} exhausts all 25 pairs of total anchor counts
by explicit cases and uses full twins where unit signatures alone
would be inadequate.
\item \textbf{Arbitrary depth.} Nonanchor crown exclusion is proved
using a support supergraph for every depth. The remaining proof
has four mode types, the terminal/nonterminal distinction, and
explicit parity inequalities. No depth cutoff is used.
\item \textbf{Shared leaves.} The strict minus threshold stops before
the leaves, both parent signs are included for leaf rows, and
the terminal unpinned case permits siblings \(p\ne q\).
\item \textbf{Dummy positions.} Dummy columns are before all variable
columns; dummy rows follow them. Their full twins and constant
nonanchor entries are used only in their stated domains.
\item \textbf{Root boundary.} At \(V_1^-\) the previous minus class
does not exist; that branch is excluded directly. For \(h=1\),
there are no nonterminal minus-horizontal cases, and the exact
count is treated separately.
\item \textbf{Repair versus hitting.} Sampling protects actual cells
of an arbitrary edited matrix. Nonroot variable pairs can be
correlated and can have arbitrary edited values. No proposed
transversal is assumed to be a free repair.
\item \textbf{Novelty and formal status.} The candidate and tree
construction are credited to prior work. No successful \(4\times4\)
theorem was located in the dated search, but publication priority
has not been proved. The new theorem has no completed Lean
formalization and has not undergone external peer review.
\end{enumerate}

Three separate mathematical derivations in this investigation agreed
on the localization argument. A further adversarial reading found two
wording overstatements in working notes: monotonicity must be restricted
to nonanchor rows, and the support graph is contained in the displayed
graph rather than necessarily equal to it. Both are corrected in this
report; neither changes the proof's required hypotheses. The exact
three-family count was independently rederived rather than inferred
from the finite table.

\section{Reproducible artifacts and executed verification}
\label{sec:verification}

\subsection{Independent implementations}
The archive contains standard-library Python programs:
\begin{itemize}
\item \texttt{code/generate.py}: emits every entry of \(A_h\) as binary
text or as an exact weighted block matrix, with deterministic ordering.
\item \texttt{code/discover.py}: independently written discovery
implementation; recursively orders the tree, assigns entries by
role predicates, inspects every row-block four-tuple, and counts
all compatible column chains with integer arithmetic.
\item \texttt{code/independent\_verify.py}: imports neither the discovery
program nor prior code. It constructs order using rational interval
endpoints and constructs signatures using explicit role sets. It
inspects every \emph{column}-block four-tuple and enumerates all
compatible row chains. It verifies the entire certificate set,
every one and zero, strict orders, exact multiplicities, and the
localized zero coordinate.
\item \texttt{code/families.py}: generates the proved three families
and compares their complete sets and integer weights with the
exhaustive certificates.
\item \texttt{code/finite\_checks.py}: checks the original and suggested
38-order anchors, all Boolean transition cases, parent-residue
inequalities, and the reversed-parity negative-control witness.
\end{itemize}
No solver output, floating-point arithmetic, or randomized search is
used for a proof claim. Timings alone use a wall clock.

\subsection{Executed complete finite searches}
\begin{center}
\small
\begin{tabular}{r r r r r r}
\toprule
\(h\)&\(n_h\)&block order&four-tuples&block copies&expanded copies\\
\midrule
1&32&17&2,380&2&128\\
2&120&35&52,360&15&23,040\\
3&352&59&455,126&78&2,080,768\\
4&928&95&3,183,545&348&152,567,808\\
5&2,304&155&23,130,030&1,464&10,385,096,704\\
\bottomrule
\end{tabular}
\end{center}
At each depth the discovery program inspected all row four-tuples;
the verifier independently inspected all column four-tuples. Every
complete weighted copy set agreed. The independent verifier completed
the five-depth run in approximately 18.75 seconds in this environment.
Complete per-embedding certificates, execution receipts, and a full
expanded depth-one host are included.

For a certificate, an anchor descriptor is
\texttt{["A", level, kind, parity, coordinate]}, with coordinates
zero-based \(0,1\), kind \texttt{V+}, \texttt{V-}, \texttt{W+}, or
\texttt{W-}, and null parity for \(V\). A variable descriptor is
\texttt{["T", level, node, sign]}, where \texttt{L} denotes the
shared leaf. Dummy is \texttt{["D"]}. Each record has four rows,
four columns, and an exact integer product weight.

Twin compression is exact: selecting two positions from one block
would repeat a full row or column of the host, impossible because
the four pattern traces on that axis are distinct. Thus a valid
compressed copy expands into exactly the product of its eight
block multiplicities, and all expanded copies arise uniquely.
This is why the finite certificates count the full host, not just
a smaller surrogate.

\subsection{Why the finite checks do not extrapolate}
The finite runs establish the stated finite copy sets and check two
independent implementations. The infinite result follows from
Lemmas~\ref{lem:structure}--\ref{lem:shape}, the exhaustive symbolic
mode cases, and the protected-path induction, all quantified over
arbitrary \(h\ge1\). Their finite case structure is independent
of \(h\): only signs, the four mode types, child residues, and the
terminal boundary are split. Parent monotonicity follows from
\(p=2\pi(p)+(p\bmod2)\); the only residue values are \(0,1\).
The exact count follows from complete symbolic families and finite
sum identities, not a fitted sequence.

\subsection{Commands}
From the extracted archive directory:
\begin{samepage}
\begin{verbatim}
python3 code/independent_verify.py --max-depth 5 \
  --output logs/recheck.json
python3 code/families.py --max-h 5
python3 code/finite_checks.py --output logs/finite_recheck.json
python3 code/generate.py --depth 2 --format txt \
  --output host_h2.txt
python3 code/discover.py --h 5 --count --collect \
  --output certificates/h4_copies_h5.json
pdflatex -interaction=nonstopmode -halt-on-error report.tex
pdflatex -interaction=nonstopmode -halt-on-error report.tex
\end{verbatim}
\end{samepage}
The first command verifies all provided certificates without modifying
them. The discovery command explicitly regenerates the named one.
The largest full binary host grows quadratically in \(n_h\) as a file;
the JSON representation preserves weights and is usually more convenient.
The programs raise explicit exceptions, so optimized Python does not
disable their checks.

\section{A kernel-checked Lean route}
\label{sec:lean}
The formalization should first target the simpler sparse-zero bound;
the exact count can be a subsequent theorem. No unproved Lean stub
is presented as an executed result here.
\begin{enumerate}
\item Define \(\Hfour:\mathrm{Fin}\ 4\to\mathrm{Fin}\ 4\to
\mathrm{Bool}\), ordered copies, and unrestricted Hamming distance
as in the existing Model module. Prove finite distinctness of its
four rows and columns with ordinary kernel-checked \texttt{decide}.
\item Parameterize the anchor-group type in the host, replacing
\(\mathrm{Fin}\ 64\) by \(\mathrm{Fin}\ 2\), and use
\(d_h=14h+2\). Reuse or adapt \texttt{TreePositions},
\texttt{HostTreeOrder}, and \texttt{HostBodyOrder}. Preserve one
literal leaf position type with both sign views.
\item Package structural hypotheses: full twins, a common anchor
prefix and mode order, one-class positive neighborhoods, nested
nonanchor traces, two signatures per variable class, and nonanchor
crown exclusion. Prove them from the concrete role and entry tables.
\item Prove the abstract anchor-count lemma for increasing maps
\(\mathrm{Fin}\ 4\to[n]\). Formalize its \((a,b)\) cases explicitly.
The \((2,2)\) branch uses full pattern traces to exclude extra
anchors; reusing a large-anchor density theorem would not suffice.
\item For \((2,1)\), extract the complementary groups and roles.
Prove the four all-depth mode cases using order equivalences and
integer division. The nonterminal minus case should expose
\(p=2b,\ q=2b+1\) and then use the concrete mode ordering.
These local arithmetic implications are suitable for
\texttt{omega} after unfolding the relevant natural-number facts.
\item Obtain an injection of copies into
\(Z_h\times[n]^3\times[n]^3\), proving
\(\#Z_h=2m\) and \(N_{\Hfour}(A_h)\le2mn^6\).
The fixed coordinate is \((3,4)\), with value zero.
\item Reuse the generic Boolean \texttt{PathPropagation} and
finite \texttt{SamplingDistance} arguments. Re-establish the
two-group frame extension and exact protected-cell sampling
bound. Quantify over every binary edited matrix.
\item Adapt \texttt{SequenceAssembly} and \texttt{Consequence}
with exponents \(8,6\) and coefficient 2. Prove the limit giving
the fully quantified failure theorem.
\item Optionally formalize the three explicit copy families and
their disjoint cardinalities to obtain the exact count.
\end{enumerate}
The final verification gate should build the concrete solution,
run the comparator on the exact new definitions and theorem, and
print its axioms. Permit only the standard logical axioms already
used upstream, with no \texttt{sorryAx}, new postulated combinatorial
lemmas, or unchecked native-computation oracle. This is a proposed
verification route, not a completed build in this investigation.

\appendix
\section{The suggested 38-by-38 anchor is valid but superseded}
\label{app:38}
For one-based \(u,v\in[38]\), start from \(J_{38}-I_{38}\).
Replace rows \(7,\ldots,38\) and columns \(34,\ldots,38\)
by the 32 lexicographically ordered five-bit words. Call this
matrix \(S_{38}\). Its first six rows and columns are unaffected
by the replacement, and every opposite trace there has at most
one zero. Each row has at least 32 ones, each column at least
22. Its rows and columns are pairwise distinct: rows \(1,\ldots,33\)
have individual zeros in the first 33 columns; the last five rows
are distinguished by different word suffixes. The first 33 columns
have individual zero positions, while each word column has 16
zeros and the five coordinate functions differ.

Use the original tree host with this anchor. Each variable position
has at most four mode incidences against distinct selected anchor
groups. Dummy twins allow at most one selected dummy on each axis.
No five nonanchor columns are shattered: if two columns lie in one
variable class, nesting forbids the two opposite mixed traces.
Otherwise five distinct classes would need a non-dummy row with
four ones and a zero to realize all words, but a variable row meets
at most three classes. A dummy column itself also prevents shattering.

If a selected \(S_{38}\)-copy has at most six anchors on each axis,
its final 32 rows and last five columns are nonanchors, contradicting
the word block. Otherwise it has at least six anchor positions on
one axis. On its first six anchors every opposite pattern trace has
at least five ones, excluding all variable positions on the other
axis. All but at most one opposite position are anchors. Two such
anchors have a common positive entry on the six-position core, so
all belong to one mode. Density excludes anchor positions of other
modes on the original axis. A remaining dummy on the opposite axis
has at most one positive entry against that mode's selected groups,
contradicting the dense core. Finally every nonanchor on the original
axis has at most one positive entry against the complete opposite
mode groups, contradicting row or column density. Full distinctness
and order exhaust all 38 groups on both axes.

Thus the entire baseline argument works and gives a \(40\times40\)
pattern. The finite anchor checks in the archive confirm the exact
properties above. This is not claimed as a novel size improvement,
because the prior \(5\times5\) construction is already much smaller.
The successful new argument does not use this anchor.

\section{Dated source scope and novelty limits}
The record distinguishes source availability from novelty.
\begin{itemize}
\item The OpenAI main commit, family-path histories, comparator,
solution implementation, public PR collection, and relevant PR
were inspected. No family-190 size change was present.
\item The ProveIt path was checked again during the investigation.
Its current main was
\texttt{1703447ff43256042649da80c73f75121c48c95b};
the OrderedMatrices path still had only the October 7 intake
commit. Its research directory contained the same \(5\times5\)
package, including the displayed \(4\times4\) candidate.
\item Searches covered the exact title, ordered binary matrix
removal, four-by-four and five-by-five variants, the relevant
repositories, and citations and follow-ups to~\cite{ab}.
No separate successful \(4\times4\) theorem was located.
\item The mathematical proof of a smaller obstruction is independent
of this absence finding. Global publication priority is not verified.
Ordinary web search initially missed the \(5\times5\) report;
the GitHub code search was materially more informative.
\end{itemize}
The accompanying audits give the reviewed URLs, dates, query scope,
and the distinction between source inspection, prior-code reruns,
and independent finite verification.

\begin{thebibliography}{9}
\raggedright
\bibitem{oai}
OpenAI, \emph{Polynomial removal fails for ordered binary matrices},
25 September 2026, family 190, snapshot
\texttt{fd4aeeb2ee4fc729c18d98444fed42fd0529eeeb}.
\url{https://github.com/openai/math/tree/fd4aeeb2ee4fc729c18d98444fed42fd0529eeeb/preprints/Polynomial-removal-fails-for-ordered-binary-matrices-September-25-2026}.
\bibitem{oaiLean}
OpenAI, family-190 Lean catalogue and MatrixRemoval implementation.
\url{https://github.com/openai/math/blob/fd4aeeb2ee4fc729c18d98444fed42fd0529eeeb/lean/docs/190.md}.
\url{https://github.com/openai/math/tree/fd4aeeb2ee4fc729c18d98444fed42fd0529eeeb/lean/OAI/Combinatorics/MatrixRemoval}.
\bibitem{proveit}
Research draft for the ProveIt project,
\emph{A \(5\times5\) Obstruction to Polynomial Ordered Matrix Removal},
7 October 2026, first path-history commit
\texttt{bfbfe915186310fdf745a7a089e7dd88eeb773b1}.
\url{https://github.com/VladimirReshetnikov/ProveIt/blob/bfbfe915186310fdf745a7a089e7dd88eeb773b1/Combinatorics/OrderedMatrices/Research/ordered-matrix-removal/ordered_matrix_removal.tex}.
The candidate and failed-repair witness are in
\texttt{sections/generalization\_and\_research.tex} at the same commit.
\bibitem{ab}
N. Alon and O. Ben-Eliezer,
\emph{Efficient Removal Lemmas for Matrices}, arXiv:1609.04235v3,
12 June 2017; \emph{Order} 37 (2020), 83--101.
\url{https://arxiv.org/abs/1609.04235}.
\bibitem{abf}
N. Alon, O. Ben-Eliezer, and E. Fischer,
\emph{Testing Hereditary Properties of Ordered Graphs and Matrices},
arXiv:1704.02367; FOCS 2017.
\url{https://arxiv.org/abs/1704.02367}.
\bibitem{gs}
L. Gishboliner and A. Shapira,
\emph{Polynomial Property Testing}, arXiv:2508.16878v1,
23 August 2025.
\url{https://arxiv.org/abs/2508.16878}.
\bibitem{gt}
L. Gishboliner and I. Tomon,
\emph{Polynomial removal lemmas for ordered graphs},
\emph{Combinatorial Theory} 2(3) (2022), Paper 3;
arXiv:2110.03577v2.
\url{https://arxiv.org/abs/2110.03577}.
\end{thebibliography}
\end{document}
